By a brick-decomposition of the plane we mean a collection $G=\left\{g_{i}\right\}$ of polyhedral disks ( $=2$-cells) such that (1) $\cup_{i=1}^{\infty} g_{i}=\mathbf{R}^{2}$, (2) if two sets $g_{i}$ and $g_{j}$ intersect, then their intersection is a broken line lying in the frontier of each of them, and (3) every point has a neighborhood which intersects at most three of the sets $g_{i}$. One way to get such a collection is to cut up the plane by horizontal lines and vertical segments so as to get a sort of "infinite brick wall," as shown in Figure 4.6.

\begin{figure}[h]
\begin{center}
  \includegraphics[alt={},max width=\textwidth]{c570eb93-f29d-4347-b75d-c251ce2386f6-034_233_796_1983_661}
\captionsetup{labelformat=empty}
\caption{Figure 4.6}
\end{center}
\end{figure}

In general, given a set $M$ in a metric space $[X, d]$, the diameter $\delta M$ of $M$ is the least upper bound of the numbers $d(P, Q)(P, Q \in M)$. (Thus $\delta M$\\
may be $\infty$.) If $G$ is a collection of subsets of $X$, then the mesh of $G$ is the least upper bound of the numbers $\delta g(g \in G)$.

Evidently we can construct a brick-decomposition $G$ of $\mathbf{R}^{2}$ with mesh as small as we please. In any case, the union of any subcollection of $G$ is a 2 -manifold with boundary. In the present case, we use bricks which are rectangular regions, with sides parallel to the edges of $\operatorname{Fr} I$, such that $\bar{I}$ is the union of a subcollection of them. And we use bricks of sufficiently small diameter so that no one of them intersects both $A_{1}$ and $A_{2}$. (See Problem 1.12.)
